The \peripheral{AFE} integrates a three-electrode potentiostat analog front end with a 12-bit dual-slope integrating ADC (DSADC) and a fully programmable bias generation block, forming a self-contained electrochemical sensing platform. The potentiostat controls the potential across a working and reference electrode while measuring the resulting cell current; the DSADC digitises the resulting voltage directly on chip. Together they are designed to support amperometric and voltammetric techniques such as chronoamperometry and cyclic voltammetry for chemical concentration measurements. All analog operating points---bias currents, DAC codes, integration time, gain---are configurable through memory-mapped registers, enabling complete software control without any external trim components.

\begin{figure}[H]
  \centering
  \includegraphics[width=0.85\linewidth]{figures/potentiostatDSBlockDiagram.png}
  \caption{AFE system block diagram. The potentiostat front end (left) conditions the electrode current; the dual-slope ADC (centre) digitises the result; the digital peripheral (right) exposes both subsystems through memory-mapped registers.}
  \label{fig:afe-sysblock}
\end{figure}

\subsubsection*{Potentiostat Front End}

The potentiostat implements the classical three-electrode electrochemical cell interface: a working electrode (WE), a reference electrode (RE), and a counter electrode (CE). A folded-cascode \emph{control amplifier} ($A_1$) forms the outer feedback loop: its non-inverting input is driven by the \emph{excitation voltage} programmed through \texttt{BIAS\_REV\_POT} (14-bit R-2R DAC), while its inverting input senses the reference electrode potential. $A_1$ drives the counter electrode such that $V_\mathrm{RE} = v_\mathrm{pattern}$, holding the WE--RE potential precisely at the programmed value. This feedback is stable across the full electrochemical bandwidth ($\leq$\,300\,Hz) with a phase margin exceeding 60\degree.

The cell current $i_\mathrm{cell}$ flowing through the counter electrode is duplicated by a \emph{regulated cascode current mirror} ($Q_1$--$Q_6$) with a unity mirroring ratio, presenting high output impedance and $R^2 > 0.999$ current replication accuracy across process corners over a measured input current range of 1\,nA to 20\,$\mu$A.

A \emph{transimpedance amplifier} (TIA, $A_2$) converts the mirrored current into a voltage
\[
  v_\mathrm{out} = v_\mathrm{cm} - i_\mathrm{out} \cdot R_f
\]
where $v_\mathrm{cm} = V_\mathrm{DD}/2$ is the analog common-mode voltage (programmed via \texttt{BIAS\_DSADC\_VCM}, midscale\,=\,8192). Referencing the output to $v_\mathrm{cm}$ centres the swing on the ADC input range so that both anodic and cathodic currents can be represented without a negative supply. The feedback resistance $R_f$ is set by \texttt{BIAS\_TIA\_G\_POT} (17-bit thermometer code): the register controls 16 series 60\,k$\Omega$ resistors plus one 1\,M$\Omega$ resistor, each switchable in parallel by an analog switch. Asserting a bit closes the corresponding switch, shorting that resistor and \emph{reducing} $R_f$. Larger codes therefore give lower gain and a wider measurable current range. The reset value 0x1FFFF closes all switches (minimum gain); writing 0x00000 gives maximum gain ($R_f \approx 2$\,M$\Omega$, current sensitivity\,$\approx$\,56\,pA).

The TIA amplifier ($A_2$) is a two-stage Miller-compensated op-amp; its tail-current bias is trimmed by \texttt{BIAS\_TC\_POT} and its Miller lead-compensation resistor by \texttt{BIAS\_LC\_POT}.

\begin{figure}[H]
  \centering
  \includegraphics[width=0.80\linewidth]{figures/potentiostat_schem.png}
  \caption{Potentiostat front-end schematic. $A_1$ (folded-cascode) is the control amplifier; $Q_1$--$Q_6$ form the regulated cascode current mirror with unity gain; $A_2$ with $R_f$ implements the transimpedance amplifier. The excitation voltage $v_{\mathrm{pattern}}$ is generated by the \texttt{BIAS\_REV\_POT} DAC.}
  \label{fig:afe-potschem}
\end{figure}

\subsubsection*{Dual-Slope ADC (DSADC)}

The DSADC is a 12-bit integrating converter. A conversion proceeds in two phases:
\begin{enumerate}
  \item \textbf{Integration}: the TIA output ($v_\mathrm{out}$) is integrated for a fixed interval $T_\mathrm{int} = \mathtt{RAMPNUM} \times T_\mathrm{clk}$ clock cycles (12-bit field \texttt{AFE\_CR}[23:12]). The internal counter increments from 0 to 4095.
  \item \textbf{Deintegration}: the integrator input is switched to a reference of opposite polarity ($V_\mathrm{ref} = V_\mathrm{DD}$). The counter decrements from 4095 toward zero. When the comparator detects the integrator output crossing the common-mode voltage $v_\mathrm{cm}$, conversion completes and the counter value is latched as the result. The deintegration time $T_\mathrm{deint} = T_\mathrm{int} \cdot V_\mathrm{in}/V_\mathrm{ref}$ is directly proportional to the input magnitude.
\end{enumerate}

\begin{figure}[H]
  \centering
  \begin{subfigure}[b]{0.52\linewidth}
    \centering
    \includegraphics[width=\linewidth]{figures/dualslopeupd.png}
    \caption{Dual-slope ADC schematic with programmable $R_1$ (sets $R_1C_1$ time constant) and $R_2$ (lead-compensation at the integration-to-deintegration transition).}
    \label{fig:afe-dsschem}
  \end{subfigure}%
  \hfill
  \begin{subfigure}[b]{0.44\linewidth}
    \centering
    \includegraphics[width=\linewidth]{figures/dualslopewave.png}
    \caption{Conversion waveform. The integrator ramps up during the fixed integration phase ($T_{\mathrm{int}}$) then ramps back to $v_{\mathrm{cm}}$ during deintegration ($T_{\mathrm{deint}} \propto V_{\mathrm{in}}$).}
    \label{fig:afe-dswave}
  \end{subfigure}
  \caption{Dual-slope ADC circuit and conversion waveform.}
  \label{fig:afe-dsadc}
\end{figure}

The 12-bit result is available in \texttt{AFE\_ADC\_VAL} bits 11:0 after \texttt{DATARDYIF} asserts. The integrating architecture provides inherent low-frequency noise rejection and high linearity; measured deviations from the ideal transfer curve are $\leq$\,1\,LSB across the full input range [0, 1.25\,V]. The worst-case sampling rate at a 6\,MHz ADC clock is 732\,Hz (when $V_\mathrm{in} = V_\mathrm{ref}$, so $T_\mathrm{deint} = T_\mathrm{int}$).

A single conversion is triggered by writing any value to \texttt{AFE\_ADC\_VAL} while \texttt{ADCEN}\,=\,1 and \texttt{CONTMEAS}\,=\,0. In continuous mode (\texttt{CONTMEAS}\,=\,1), the FSM restarts immediately after each conversion without software intervention.

\begin{figure}[H]
  \centering
  \begin{subfigure}[b]{0.44\linewidth}
    \centering
    \includegraphics[width=\linewidth]{figures/fsmstatediagram.png}
    \caption{DSADC FSM state diagram. The FSM exits idle on a write to \texttt{AFE\_ADC\_VAL} (or immediately in continuous mode), runs the integration and deintegration phases, then asserts \texttt{DATARDYIF}.}
    \label{fig:afe-fsm}
  \end{subfigure}%
  \hfill
  \begin{subfigure}[b]{0.52\linewidth}
    \centering
    \includegraphics[width=\linewidth]{figures/dsTimingDia.png}
    \caption{DSADC control signal timing diagram showing the reset, integration, and deintegration phases together with the counter direction and the conversion-done pulse.}
    \label{fig:afe-dstiming}
  \end{subfigure}
  \caption{DSADC finite state machine and timing signals.}
  \label{fig:afe-fsm-timing}
\end{figure}

The DSADC input is selected by \texttt{ADCEXTIN} (\texttt{AFE\_CR} bit\,8): 1\,=\,potentiostat output pad, 0\,=\,external pin. The integrating amplifier (also a two-stage BTS op-amp) has its tail-current bias trimmed by \texttt{BIAS\_TC\_DSADC} and its Miller resistor by \texttt{BIAS\_LC\_DSADC}. The DSADC integration resistor $R_1$ (which sets the $R_1C_1$ time constant) is trimmed by \texttt{BIAS\_RIN\_DSADC}, and the lead-compensation resistor $R_2$ (which mitigates the integrator overshoot at the integration-to-deintegration transition) is trimmed by \texttt{BIAS\_RFB\_DSADC}. Both on-chip resistors can vary $\pm$20\,\% from nominal, and these 6-bit trim codes correct for process variation.

\subsubsection*{Bias Generation}

Both the potentiostat and DSADC require supply-independent bias voltages and currents. The source is selected by \texttt{BIAS\_USEDAC} (\texttt{BIAS\_CR} bit\,4):

\begin{itemize}
    \item \textbf{Internal bias generator} (\texttt{BIAS\_USEDAC}\,=\,0): a wide-swing cascode generator produces the global bias voltages $V_{bp}$, $V_{bpc}$, $V_{bnc}$, $V_{bn}$ distributed to all analog blocks. Enable with \texttt{BIAS\_EN} (\texttt{BIAS\_CR} bit\,2). The overall current level is trimmed by \texttt{BIAS\_ADJ} (6-bit; higher codes produce smaller currents). Enable output buffers with \texttt{BIAS\_BUFEN} (\texttt{BIAS\_CR} bit\,3).
    \item \textbf{External bias DACs} (\texttt{BIAS\_USEDAC}\,=\,1): four 14-bit R-2R DACs (\texttt{BIAS\_DBP}, \texttt{BIAS\_DBPC}, \texttt{BIAS\_DBNC}, \texttt{BIAS\_DBN}) set the P, P-cascode, N-cascode, and N bias voltages independently, allowing fine-grained tuning or corner compensation without relying on the on-chip generator.
\end{itemize}

Two additional 14-bit R-2R DACs generate voltages that must be enabled via \texttt{DACEN} (\texttt{AFE\_CR} bit\,3) before a conversion is started:
\begin{itemize}
    \item \texttt{BIAS\_REV\_POT}: generates the excitation voltage $v_\mathrm{pattern}$ applied to the non-inverting input of the control amplifier. Programming this register sets the WE--RE electrode potential. For cyclic voltammetry, firmware sweeps this register over time; for chronoamperometry, it is stepped to a fixed overpotential.
    \item \texttt{BIAS\_DSADC\_VCM}: generates the analog common-mode voltage $v_\mathrm{cm} = V_\mathrm{DD}/2$ applied to both $A_1$ and $A_2$. Midscale code 8192 produces $V_\mathrm{DD}/2$, centring the TIA output and DSADC input swing.
\end{itemize}

\subsubsection*{Interrupts and Flags}

Two status flags in \texttt{AFE\_SR} indicate conversion state:

\begin{itemize}
    \item \textbf{DATARDYIF} (SR bit\,1): set by hardware when a DSADC conversion completes and the result is latched in \texttt{AFE\_ADC\_VAL}. Write\,1 to clear. An IRQ is generated if \texttt{DATARDYIE} (\texttt{AFE\_CR} bit\,2) is also set.
    \item \textbf{OVFIF} (SR bit\,2): set when a conversion completes while \texttt{DATARDYIF} is still asserted, indicating the previous result was overwritten before it was read. Write\,1 to clear.
\end{itemize}

\subsubsection*{Operation}

\begin{figure}[H]
  \centering
  \includegraphics[width=0.75\linewidth]{figures/peripheral.png}
  \caption{AFE digital peripheral block. The memory-mapped register interface connects the microcontroller bus to the DSADC FSM, bias DACs, and test port multiplexer. Clock gating minimises power when the AFE is idle.}
  \label{fig:afe-peripheral}
\end{figure}

\begin{enumerate}
    \item Enable and configure bias generation: set \texttt{BIAS\_EN} and \texttt{BIAS\_BUFEN} in \texttt{BIAS\_CR}. Write \texttt{BIAS\_ADJ} (nominal code is chip-dependent; see characterization data), or configure the individual bias DACs and set \texttt{BIAS\_USEDAC}\,=\,1.
    \item Write the TIA and DSADC amplifier bias trim codes: \texttt{BIAS\_TC\_POT}, \texttt{BIAS\_LC\_POT}, \texttt{BIAS\_TC\_DSADC}, \texttt{BIAS\_LC\_DSADC}.
    \item Write the DSADC resistor trim codes: \texttt{BIAS\_RIN\_DSADC} (integration resistor $R_1$) and \texttt{BIAS\_RFB\_DSADC} (lead-compensation resistor $R_2$).
    \item Set the TIA gain by writing \texttt{BIAS\_TIA\_G\_POT} (0x1FFFF\,=\,minimum gain / maximum current range; 0x00000\,=\,maximum gain / minimum current range).
    \item Program the electrode potential: write \texttt{BIAS\_REV\_POT} to set $v_\mathrm{pattern}$ (i.e.\ the WE--RE voltage).
    \item Program the analog common mode: write \texttt{BIAS\_DSADC\_VCM}\,=\,8192 for $V_\mathrm{DD}/2$.
    \item Enable the DACs by setting \texttt{DACEN}\,=\,1 in \texttt{AFE\_CR}.
    \item Set the integration length: write \texttt{RAMPNUM} (\texttt{AFE\_CR} bits 23:12). The integration time is $\mathtt{RAMPNUM} \times T_\mathrm{clk}$.
    \item Enable the AFE clock and ADC: set \texttt{EN}\,=\,1 and \texttt{ADCEN}\,=\,1 in \texttt{AFE\_CR}.
    \item Trigger a conversion by writing any value to \texttt{AFE\_ADC\_VAL} (single-shot), or set \texttt{CONTMEAS}\,=\,1 for continuous mode.
    \item Wait for \texttt{DATARDYIF} (poll \texttt{AFE\_SR} bit\,1 or wait for the IRQ).
    \item Read the 12-bit result from \texttt{AFE\_ADC\_VAL} bits 11:0. Clear \texttt{DATARDYIF} by writing\,1 to \texttt{AFE\_SR} bit\,1.
\end{enumerate}

For \textbf{cyclic voltammetry}, repeat steps\,10--12 while periodically updating \texttt{BIAS\_REV\_POT} to sweep the electrode potential over the desired voltage window. For \textbf{chronoamperometry}, hold \texttt{BIAS\_REV\_POT} at the step potential and use continuous mode (\texttt{CONTMEAS}\,=\,1) to sample the decaying cell current as a function of time.

\subsubsection*{Digital Test Ports}

The AFE provides four digital test ports (\texttt{DTP0}--\texttt{DTP3}), each driven by a selected signal from a 32-entry internal debug vector. The 5-bit fields \texttt{DTP0SEL} through \texttt{DTP3SEL} in \texttt{AFE\_TPR} select the signal for each port. All four ports are always configured as outputs. The available signals are listed in Table~\ref{t:afe-dtp}.

\begin{table}[h]
\centering
\begin{tabularx}{\textwidth}{c l X}
\hline
\textbf{Index} & \textbf{Signal} & \textbf{Description} \\
\hline
0  & \texttt{ADCACTIVE}    & High while the DSADC FSM is running a conversion (\texttt{AFE\_SR} bit~0). \\
1  & \texttt{DATARDYIF}    & Data-ready interrupt flag (\texttt{AFE\_SR} bit~1). High when a conversion result is available. \\
2  & \texttt{OVFIF}        & Overflow flag (\texttt{AFE\_SR} bit~2). High when a conversion completed before the previous result was read. \\
3  & \texttt{adc\_done}    & Internal done pulse from the AFE FSM. High for one clock cycle when a conversion completes. \\
4  & \texttt{adc\_clk}     & DSADC conversion clock output, derived from the AFE clock. \\
5  & \texttt{adc\_switch[0]} & DSADC switch control bit 0 from the AFE FSM. \\
6  & \texttt{adc\_switch[1]} & DSADC switch control bit 1. \\
7  & \texttt{adc\_switch[2]} & DSADC switch control bit 2. \\
8  & \texttt{adc\_start}   & ADC start signal. High continuously in continuous mode; single pulse when \texttt{AFE\_ADC\_VAL} is written in single-shot mode. \\
9--19  & ---               & Tied low. Not connected to any internal signal. \\
20 & \texttt{DSADC\_CNT[0]}  & Bit~0 (LSB) of the live DSADC ramp counter. \\
21 & \texttt{DSADC\_CNT[1]}  & Bit~1 of the DSADC ramp counter. \\
22 & \texttt{DSADC\_CNT[2]}  & Bit~2. \\
23 & \texttt{DSADC\_CNT[3]}  & Bit~3. \\
24 & \texttt{DSADC\_CNT[4]}  & Bit~4. \\
25 & \texttt{DSADC\_CNT[5]}  & Bit~5. \\
26 & \texttt{DSADC\_CNT[6]}  & Bit~6. \\
27 & \texttt{DSADC\_CNT[7]}  & Bit~7. \\
28 & \texttt{DSADC\_CNT[8]}  & Bit~8. \\
29 & \texttt{DSADC\_CNT[9]}  & Bit~9. \\
30 & \texttt{DSADC\_CNT[10]} & Bit~10. \\
31 & \texttt{DSADC\_CNT[11]} & Bit~11 (MSB) of the live DSADC ramp counter. Observing bits 31:20 simultaneously gives the full 12-bit live counter value. \\
\hline
\end{tabularx}
\caption{AFE digital test port debug signal index map.}
\label{t:afe-dtp}
\end{table}
